\documentclass[12pt]{article}
\usepackage[utf8]{inputenc}
\usepackage[T1]{fontenc}
\usepackage[margin=1in]{geometry}
\usepackage{setspace}
\usepackage{xcolor}
\usepackage[colorlinks=true, urlcolor=blue!60!black, linkcolor=blue!60!black]{hyperref}
\setlength{\parindent}{2em}
\setlength{\parskip}{1.25em}
\renewcommand{\baselinestretch}{1.0}
\begin{document}
\begin{center}
{\LARGE \textit{Romer v. Evans}, 517 U.S. 620 (1996)\par}
\vspace{0.8em}
{\large Case Brief\par}
\end{center}
\vspace{1em}

\section{Memory Jogger}

In \href{https://www.courtlistener.com/opinion/118027/romer-v-evans/}{\textit{Romer v. Evans}}, the Supreme Court struck down Colorado's Amendment 2 under the Fourteenth Amendment's Equal Protection Clause, holding that a state constitutional ban forbidding all state and local anti-discrimination protections for gay, lesbian, and bisexual people fails even rational-basis review because its sweeping disability is inexplicable by anything other than animus toward the targeted class.

\section{Detailed Case Facts}

Prior to 1992, several Colorado municipalities---including Aspen (Aspen Mun. Code \S~13-98 (1977)), Boulder (Boulder Rev. Code \S\S~12-1-1 to 12-1-11 (1987)), and the City and County of Denver (Denver Rev. Mun. Code, Art. IV, \S\S~28-91 to 28-116 (1991))---enacted ordinances banning discrimination on the basis of sexual orientation in housing, employment, education, public accommodations, and health and welfare services (\href{https://supreme.justia.com/cases/federal/us/517/620/}{517 U.S. at 623--24}). State-level measures likewise barred sexual-orientation discrimination, including Governor Roy Romer's Executive Order No. D0035 (Dec. 10, 1990) covering state employees, Colorado Insurance Code \S~10-3-1104 barring insurers from basing insurability or premiums on sexual orientation, and nondiscrimination policies at state colleges such as Metropolitan State College of Denver and Colorado State University (\href{https://supreme.justia.com/cases/federal/us/517/620/}{517 U.S. at 626--27}).

In November 1992, following a contentious statewide campaign, Colorado voters approved a ballot initiative---designated ``Amendment 2'' and codified as Article II, \S~30b of the Colorado Constitution---by a margin of approximately 53\% to 47\% (\href{https://supreme.justia.com/cases/federal/us/517/620/}{517 U.S. at 623--24}; \href{https://en.wikipedia.org/wiki/Romer_v._Evans}{Wikipedia}). Entitled ``No Protected Status Based on Homosexual, Lesbian or Bisexual Orientation,'' Amendment 2 provided:
\begin{quote}
``Neither the State of Colorado, through any of its branches or departments, nor any of its agencies, political subdivisions, municipalities or school districts, shall enact, adopt or enforce any statute, regulation, ordinance or policy whereby homosexual, lesbian or bisexual orientation, conduct, practices or relationships shall constitute or otherwise be the basis of or entitle any person or class of persons to have or claim any minority status, quota preferences, protected status or claim of discrimination. This Section of the Constitution shall be in all respects self-executing.'' (\href{https://supreme.justia.com/cases/federal/us/517/620/}{517 U.S. at 624}.)
\end{quote}

Immediately after Amendment 2's passage, Richard G. Evans (a gay man employed by the City of Denver), other gay and lesbian individuals, and the cities of Aspen, Boulder, and Denver sued in state court to declare Amendment 2 unconstitutional and enjoin its enforcement, naming Governor Roy Romer (who had personally opposed Amendment 2), the Colorado Attorney General, and the State of Colorado as official-capacity defendants (\href{https://supreme.justia.com/cases/federal/us/517/620/}{517 U.S. at 625}; \href{https://www.oyez.org/cases/1995/94-1039}{Oyez}).

\section{Procedural History}

\begin{enumerate}
\item \textbf{Colorado District Court (Preliminary Injunction, 1993):} The District Court for the City and County of Denver granted a preliminary injunction staying enforcement of Amendment 2 (\href{https://supreme.justia.com/cases/federal/us/517/620/}{517 U.S. at 625}).
\item \textbf{Colorado Supreme Court (\textit{Evans I}, 854 P.2d 1270 (Colo. 1993)):} Sustaining the preliminary injunction and remanding for trial, the Colorado Supreme Court held that Amendment 2 was subject to strict scrutiny under the Equal Protection Clause because it infringed the fundamental right of an independently identifiable class to participate equally in the political process, relying on voting-rights and political-restructuring precedents such as \textit{Hunter v. Erickson} and \textit{Washington v. Seattle School District No. 1} (\href{https://supreme.justia.com/cases/federal/us/517/620/}{517 U.S. at 625}). The U.S. Supreme Court denied interlocutory certiorari (510 U.S. 959 (1993)).
\item \textbf{Colorado District Court (Permanent Injunction, 1993):} On remand, the trial court rejected plaintiffs' claim that gays and lesbians constitute a suspect or quasi-suspect class (a ruling plaintiffs did not appeal), but held that the State's asserted interests failed to satisfy strict scrutiny and permanently enjoined Amendment 2 (\href{https://supreme.justia.com/cases/federal/us/517/620/}{517 U.S. at 625--26, 640 n.1}).
\item \textbf{Colorado Supreme Court (\textit{Evans II}, 882 P.2d 1335 (Colo. 1994)):} The Colorado Supreme Court affirmed the permanent injunction in a second opinion (\href{https://supreme.justia.com/cases/federal/us/517/620/}{517 U.S. at 626}).
\item \textbf{U.S. Supreme Court (517 U.S. 620 (1996)):} The Supreme Court granted certiorari (513 U.S. 1146 (1995)), heard oral argument on October 10, 1995, and on May 20, 1996, affirmed the Colorado Supreme Court's judgment on a different rationale---rational-basis review under the Equal Protection Clause (\href{https://supreme.justia.com/cases/federal/us/517/620/}{517 U.S. at 620, 626}; \href{https://www.oyez.org/cases/1995/94-1039}{Oyez}).
\end{enumerate}

\section{Judicial Votes}

The decision was \textbf{6--3} affirming the invalidation of Colorado's Amendment 2 (\href{https://supreme.justia.com/cases/federal/us/517/620/}{517 U.S. at 621}):
\begin{itemize}
\item \textbf{Majority (6):} Justice Anthony M. Kennedy (author of the opinion of the Court), joined by Justices John Paul Stevens, Sandra Day O'Connor, David H. Souter, Ruth Bader Ginsburg, and Stephen G. Breyer.
\item \textbf{Concurring Opinions (0):} None.
\item \textbf{Dissenting Opinion (3):} Justice Antonin Scalia, joined by Chief Justice William H. Rehnquist and Justice Clarence Thomas.
\end{itemize}

\section{Holding}

\textbf{Judgment:} Affirmed. Amendment 2 of the Colorado Constitution violates the Equal Protection Clause of the Fourteenth Amendment (\href{https://supreme.justia.com/cases/federal/us/517/620/}{517 U.S. at 635--36}).

\textbf{Rule Announced:} Even when a challenged classification neither burdens a recognized fundamental right nor targets a suspect or quasi-suspect class, it must bear a rational relationship to an independent and legitimate governmental purpose. A state constitutional enactment that identifies a single class of citizens by one trait and imposes a broad, undifferentiated disability across the board---barring them from seeking specific legal protection against discrimination at every level of state and local government---defies conventional rational-basis review, raises the inevitable inference of animosity toward the affected class, and violates the Equal Protection Clause because a ``bare \dots{} desire to harm a politically unpopular group'' cannot constitute a legitimate governmental interest (\href{https://supreme.justia.com/cases/federal/us/517/620/}{517 U.S. at 631--35}).

\textbf{Holding vs. Dicta:} The Court's holding is that Amendment 2 fails rational-basis scrutiny under the Equal Protection Clause because it imposes an unprecedented, sweeping legal disability on gay, lesbian, and bisexual citizens that is divorced from any legitimate, discrete state objective. Because the Court invalidated Amendment 2 on rational-basis grounds, it expressly declined to adopt the Colorado Supreme Court's strict-scrutiny ``fundamental right to participate in the political process'' rationale, did not decide whether sexual orientation is a suspect or quasi-suspect class, and did not decide whether Amendment 2 also stripped gays and lesbians of generally applicable protections against arbitrary governmental or private discrimination (though Part II suggested that was a ``fair, if not necessary, inference,'' which operated as persuasive dicta) (\href{https://supreme.justia.com/cases/federal/us/517/620/}{517 U.S. at 626, 630--31}). The majority also did not mention or overrule \textit{Bowers v. Hardwick}, 478 U.S. 186 (1986), leaving that tension until \textit{Lawrence v. Texas}, 539 U.S. 558 (2003).

\section{Analysis of Opinions}

\subsection{Majority Opinion (Justice Kennedy)}

Opening with Justice Harlan's dissent in \textit{Plessy v. Ferguson}---that the Constitution ``neither knows nor tolerates classes among citizens''---Justice Kennedy first dismantled Colorado's primary defense that Amendment 2 merely denied homosexuals ``special rights'' (\href{https://supreme.justia.com/cases/federal/us/517/620/}{517 U.S. at 623, 626}). Relying on the Colorado Supreme Court's authoritative construction, Kennedy showed that Amendment 2 immediately repealed municipal ordinances in Aspen, Boulder, and Denver as well as state executive orders, insurance rules, and university policies, and permanently barred every branch and subdivision of state and local government from adopting similar protections unless the state constitution were amended (\href{https://supreme.justia.com/cases/federal/us/517/620/}{517 U.S. at 626--30}). Surveying modern public-accommodations and civil-rights statutes, Kennedy explained that enumerating protected traits (such as age, marital status, military status, disability, and sexual orientation) is the standard legal device for making anti-discrimination duties concrete. Far from denying ``special rights,'' Amendment 2 imposed a ``special disability upon those persons alone,'' withholding ordinary civic safeguards ``taken for granted by most people either because they already have them or do not need them'' (\href{https://supreme.justia.com/cases/federal/us/517/620/}{517 U.S. at 628--31}).

Applying equal protection doctrine, Kennedy noted that ordinary laws classifying without a suspect class or fundamental right are upheld if rationally related to a legitimate end, as in \textit{New Orleans v. Dukes}, \textit{Williamson v. Lee Optical}, and \textit{Railway Express Agency v. New York}, because those narrow, context-bound statutes allowed courts to discern a link between classification and legitimate purpose (\href{https://supreme.justia.com/cases/federal/us/517/620/}{517 U.S. at 631--33}). Amendment 2, however, ``fails, indeed defies, even this conventional inquiry'' for two interrelated reasons (\href{https://supreme.justia.com/cases/federal/us/517/620/}{517 U.S. at 632}):
\begin{enumerate}
\item \textbf{Unprecedented Structural Disqualification:} Amendment 2 is ``at once too narrow and too broad,'' identifying people by a single trait and disqualifying them across the board from seeking aid from government---a ``denial of equal protection of the laws in the most literal sense'' (\href{https://supreme.justia.com/cases/federal/us/517/620/}{517 U.S. at 633}). Kennedy rejected the dissent's reliance on \textit{Davis v. Beason}, 133 U.S. 333 (1890), noting that its advocacy holding was overruled by \textit{Brandenburg v. Ohio}, any status-based disenfranchisement would fail strict scrutiny, and felon disenfranchisement was not implicated (\href{https://supreme.justia.com/cases/federal/us/517/620/}{517 U.S. at 634}).
\item \textbf{Inference of Impermissible Animus:} Quoting \textit{Department of Agriculture v. Moreno}, 413 U.S. 528, 534 (1973), Kennedy held that Amendment 2's sheer breadth was so discontinuous with Colorado's asserted rationales---protecting landlords' and employers' freedom of association and conserving enforcement resources for other groups---that it raised ``the inevitable inference that the disadvantage imposed is born of animosity toward the class of persons affected'' (\href{https://supreme.justia.com/cases/federal/us/517/620/}{517 U.S. at 634--35}). Concluding that Amendment 2 was ``a classification of persons undertaken for its own sake'' that deemed a class ``a stranger to its laws,'' the Court invalidated it (\href{https://supreme.justia.com/cases/federal/us/517/620/}{517 U.S. at 635}).
\end{enumerate}

\subsection{Dissenting Opinion (Justice Scalia, joined by Chief Justice Rehnquist and Justice Thomas)}

Justice Scalia opened with a memorable rebuke: ``The Court has mistaken a \textit{Kulturkampf} for a fit of spite'' (\href{https://supreme.justia.com/cases/federal/us/517/620/}{517 U.S. at 636}). He mounted four principal arguments:
\begin{enumerate}
\item \textbf{No Denial of Equal Treatment:} Citing footnote 9 of \textit{Evans II}, Scalia argued that the Colorado Supreme Court had already confirmed Amendment 2 left generally applicable anti-discrimination laws intact and only barred ``special treatment'' for homosexuality. Treating a requirement that a group seek preferential protections at the statewide constitutional level rather than the municipal level as an equal protection violation was ``terminal silliness'' that would destabilize multilevel democracy (\href{https://supreme.justia.com/cases/federal/us/517/620/}{517 U.S. at 637--40}).
\item \textbf{The \textit{Bowers v. Hardwick} Syllogism:} Highlighting the majority's complete silence on \textit{Bowers v. Hardwick}, 478 U.S. 186 (1986) (which respondents had not asked to overrule), Scalia argued that if a state may constitutionally criminalize homosexual conduct, then \textit{a fortiori} it is rational for a state to enact the lesser measure of prohibiting special anti-discrimination protections for such conduct or for orientation as a proxy for conduct (\href{https://supreme.justia.com/cases/federal/us/517/620/}{517 U.S. at 640--43}). Under \textit{United States v. Salerno}, a facial challenge had to fail because Amendment 2 was plainly valid as applied to conduct (\href{https://supreme.justia.com/cases/federal/us/517/620/}{517 U.S. at 643}).
\item \textbf{Moral Disapproval and Polygamy Precedent:} Scalia contended that Coloradans---who had repealed their anti-sodomy statute in 1971---were entitled to preserve traditional sexual morality through a statewide democratic referendum countering the geographic concentration and political influence of gay-rights advocates in cities like Aspen, Boulder, and Denver (\href{https://supreme.justia.com/cases/federal/us/517/620/}{517 U.S. at 644--47}). He pointed to state constitutional bans on polygamy in Arizona, Idaho, New Mexico, Oklahoma, and Utah (required by Congress as conditions of statehood) and to \textit{Davis v. Beason} as proof that statewide constitutional rules targeting disapproved sexual practices were squarely within American constitutional tradition (\href{https://supreme.justia.com/cases/federal/us/517/620/}{517 U.S. at 648--51}).
\item \textbf{Judicial Elitism:} Scalia accused the majority of taking sides in a cultural debate on behalf of ``the lawyer class'' (citing the Association of American Law Schools' nondiscrimination policy for law-school recruiters), declaring that striking down Amendment 2 was ``an act, not of judicial judgment, but of political will'' (\href{https://supreme.justia.com/cases/federal/us/517/620/}{517 U.S. at 652--53}).
\end{enumerate}

\section{Examples: Future Applications}

\begin{enumerate}
\item \textbf{Same-Side (Invalidated): Statewide Constitutional Ban on Protections for Unmarried Cohabitants.} Voters adopt a state constitutional amendment barring the state and all municipalities from enacting or enforcing any statute, ordinance, or policy prohibiting discrimination in housing, employment, insurance, or public services based on ``unmarried cohabitation status.'' \textit{Why:} Like Amendment 2, this measure singles out a non-suspect class by one trait and imposes a sweeping, across-the-board disability barring them from ordinary political recourse at every level of government, defying any discrete legitimate rationale and triggering \textit{Romer}'s anti-animus rule (\href{https://supreme.justia.com/cases/federal/us/517/620/}{517 U.S. at 632--35}).
\item \textbf{Same-Side (Invalidated): Federal Statute Denying Recognition to Lawfully Performed Same-Sex Marriages.} Congress enacts a statute stripping federal marital benefits and protections across more than 1,000 federal laws from same-sex couples validly married under state law, while leaving all other state-recognized marriages intact. \textit{Why:} The statute imposes an unusually broad, status-based disability on a single class across federal law, raising the inference of impermissible animus under \textit{Romer} (precisely the reasoning Justice Kennedy later invoked in \href{https://www.law.cornell.edu/constitution-conan/amendment-14/section-1/sexual-orientation}{\textit{United States v. Windsor}}, 570 U.S. 744 (2013)).
\item \textbf{Opposite-Side (Upheld): State Preemption Statute Setting Uniform Statewide Employment Standards.} A state legislature passes a statute establishing a uniform statewide employment code that preempts all municipal employment ordinances---including local minimum-wage, paid-leave, and local protected-class rules---so employers face a single regulatory regime statewide. \textit{Why:} Unlike Amendment 2, this law does not single out one disfavored group by a single trait to strip it alone of legal protection; it is a facially neutral allocation of regulatory authority rationally tied to an identifiable, legitimate interest in statewide commercial uniformity (\href{https://supreme.justia.com/cases/federal/us/517/620/}{517 U.S. at 632--33}).
\item \textbf{Opposite-Side (Upheld): Narrow Religious-Exemption Clause in a Fair Housing Act.} A state enacts a comprehensive anti-discrimination statute covering sexual orientation but exempts owner-occupied duplexes of two units or fewer and religious organizations from housing-nondiscrimination mandates. \textit{Why:} Unlike Amendment 2's sweeping, undifferentiated disability, this exemption is narrow in scope and grounded in a specific factual context directly tied to legitimate interests in intimate association and religious liberty (\href{https://supreme.justia.com/cases/federal/us/517/620/}{517 U.S. at 632--33, 635}).
\item \textbf{Fence-Sitter (Borderline): Municipal Charter Amendment Repealing a Single City's Sexual-Orientation Ordinance.} Voters in a single city (such as Cincinnati's 1993 Ballot Issue 3) amend the city charter to repeal the city's sexual-orientation anti-discrimination ordinance and forbid the city council from enacting future sexual-orientation protections, while leaving state-level remedies unaffected (\href{https://en.wikipedia.org/wiki/Romer_v._Evans}{Wikipedia}). \textit{Why:} This is a classic borderline case. Challengers argue that it mirrors Amendment 2 by singling out one group for a structural political disability motivated by animus; defenders argue that because it operates only at the local level rather than preempting an entire state's legal system, voters are merely choosing not to enact a local ordinance and affected citizens remain free to seek statewide legislative protection. (Indeed, the Sixth Circuit upheld Cincinnati's charter amendment after \textit{Romer} in \textit{Equality Foundation of Greater Cincinnati, Inc. v. City of Cincinnati}, 128 F.3d 289 (6th Cir. 1997).)
\end{enumerate}

\section{Critique}

\textbf{Originalist and Textualist Critiques:} Originalist and conservative scholars, echoing Justice Scalia's dissent, have criticized \textit{Romer} as doctrinally unmoored and anti-democratic. First, critics emphasize the majority's glaring failure to confront or distinguish \textit{Bowers v. Hardwick}---a gap Professor Lynn Baker famously called \href{https://scholar.law.colorado.edu/cgi/viewcontent.cgi?article=1890&context=lawreview}{\textit{The Missing Pages of the Majority Opinion in Romer v. Evans}}, 68 U. Colo. L. Rev. 387 (1997). Second, textualists argue that the Court distorted deferential rational-basis review into an unannounced ``rational basis with bite'' without articulating objective standards for distinguishing legitimate moral legislation from unconstitutional ``animus'' (\href{https://scholarship.law.duke.edu/cgi/viewcontent.cgi?article=1081&context=dlj}{Duke Law Journal}; \href{https://scholarship.law.wm.edu/cgi/viewcontent.cgi?article=1450&context=wmborj}{William \& Mary Bill of Rights Journal}).

\textbf{Progressive and Doctrinal Perspectives:} Progressive and mainstream constitutional scholars have been sharply divided---praising \textit{Romer}'s outcome while debating its opaque reasoning. On the critical side, Evan Gerstmann warned that \textit{Romer}'s ``sloppy reasoning'' left equal protection law ``even murkier than before'' because the Court declined to recognize sexual orientation as a suspect or quasi-suspect class and instead relied on ad hoc judicial sentiment (\href{https://en.wikipedia.org/wiki/Romer_v._Evans}{Wikipedia}; \href{https://scholarlycommons.law.emory.edu/context/faculty-articles/article/1144/viewcontent/HutchinsonNotPolPowALR2014.pdf}{Hutchinson, Emory Law}). Conversely, several leading theorists offered principled reconstructions of \textit{Romer}: Daniel Farber and Suzanna Sherry defended it through \href{https://scholarship.law.vanderbilt.edu/cgi/viewcontent.cgi?article=1392&context=faculty-publications}{\textit{The Pariah Principle}}, arguing that equal protection forbids the state from branding any group---suspect or not---as social outcasts; Akhil Reed Amar analogized Amendment 2 to an unconstitutional bill of attainder in \href{https://repository.law.umich.edu/cgi/viewcontent.cgi?params=/context/mlr/article/4570}{\textit{Attainder and Amendment 2: Romer's Rightness}}, 95 Mich. L. Rev. 203 (1996); and Cass Sunstein praised \textit{Romer} as a model of judicial minimalism that invalidated a uniquely sweeping enactment without prematurely resolving broader questions like same-sex marriage (\href{http://scholarship.law.cornell.edu/cgi/viewcontent.cgi?article=1228&context=cjlpp}{Sunstein, 110 Harv. L. Rev. 4 (1996)}; \href{https://harvardlawreview.org/print/vol-136/emromer-em-has-it/}{\textit{Romer Has It}, 136 Harv. L. Rev. (2023)}).

\textbf{Underlying Competing Values:} \textit{Romer} pits three enduring constitutional values against one another: (1) the anti-caste principle and equal civic standing---ensuring no group is declared a legal ``stranger'' barred from ordinary political recourse; (2) majoritarian democracy and popular sovereignty---specifically, the power of state voters to settle contentious social issues by statewide referendum and to protect traditional moral norms or associative freedom; and (3) doctrinal clarity and judicial restraint in equal protection review.

\section{Key Quotations}

\begin{quote}
``One century ago, the first Justice Harlan admonished this Court that the Constitution `neither knows nor tolerates classes among citizens.' \textit{Plessy v. Ferguson}, 163 U.S. 537, 559 (1896) (dissenting opinion). Unheeded then, those words now are understood to state a commitment to the law's neutrality where the rights of persons are at stake.''
\end{quote}
\noindent --- Justice Kennedy, Opinion of the Court, \href{https://supreme.justia.com/cases/federal/us/517/620/}{517 U.S. at 623} (\href{https://www.courtlistener.com/opinion/118027/romer-v-evans/}{CourtListener}).

\begin{quote}
``A law declaring that in general it shall be more difficult for one group of citizens than for all others to seek aid from the government is itself a denial of equal protection of the laws in the most literal sense.''
\end{quote}
\noindent --- Justice Kennedy, Opinion of the Court, \href{https://supreme.justia.com/cases/federal/us/517/620/}{517 U.S. at 633} (\href{https://www.courtlistener.com/opinion/118027/romer-v-evans/}{CourtListener}).

\begin{quote}
``[I]f the constitutional conception of `equal protection of the laws' means anything, it must at the very least mean that a bare \dots{} desire to harm a politically unpopular group cannot constitute a legitimate governmental interest.''
\end{quote}
\noindent --- Justice Kennedy, Opinion of the Court (quoting \textit{Department of Agriculture v. Moreno}, 413 U.S. 528, 534 (1973)), \href{https://supreme.justia.com/cases/federal/us/517/620/}{517 U.S. at 634--35} (\href{https://www.courtlistener.com/opinion/118027/romer-v-evans/}{CourtListener}).

\begin{quote}
``We must conclude that Amendment 2 classifies homosexuals not to further a proper legislative end but to make them unequal to everyone else. This Colorado cannot do. A State cannot so deem a class of persons a stranger to its laws.''
\end{quote}
\noindent --- Justice Kennedy, Opinion of the Court, \href{https://supreme.justia.com/cases/federal/us/517/620/}{517 U.S. at 635} (\href{https://www.courtlistener.com/opinion/118027/romer-v-evans/}{CourtListener}).

\begin{quote}
``The Court has mistaken a \textit{Kulturkampf} for a fit of spite. The constitutional amendment before us here is not the manifestation of a ``bare \dots{} desire to harm'' homosexuals, \textit{ante}, at 634, but is rather a modest attempt by seemingly tolerant Coloradans to preserve traditional sexual mores against the efforts of a politically powerful minority to revise those mores through use of the laws.''
\end{quote}
\noindent --- Justice Scalia, dissenting (joined by Chief Justice Rehnquist and Justice Thomas), \href{https://supreme.justia.com/cases/federal/us/517/620/}{517 U.S. at 636} (\href{https://www.courtlistener.com/opinion/118027/romer-v-evans/}{CourtListener}).

\end{document}
